%% kapa: exam=1 round=22 year=2011 period=1 form=A subject=economics subject_ko=경제원론 seq=22 number=62
%% source: Q-Net 감정평가사 기출문제 (cst00309) artlSeq=5036790 file=extracted/1008/5036790/2011년 제22회 감정평가사 1차시험 기출문제/1교시.hwp
%% conversion: kordoc 4.24.51 -> tools/gen1.py (auto); review=proofread-sonnet; flags=text-verified-hwp
\begin{question}{62}
수요독점 노동시장에 관한 설명으로 옳지 않은 것은? [단, 노동공급곡선은 우상향, 노동의 한계수입생산(marginal revenue product)곡선은 우하향, 이윤을 극대화하는 수요독점기업은 상품시장에서도 독점기업임]
\begin{choices}
  \item 이 노동시장의 균형고용량은 완전경쟁 노동시장의 균형고용량보다 적다.
  \item 이 노동시장의 균형임금과 완전경쟁 노동시장의 균형임금 사이에 최저임금을 강제적으로 설정할 경우 고용량이 증가할 수 있다.
  \item 이 노동시장의 균형임금은 노동의 한계수입생산보다 낮은 수준에서 결정된다.
  \item 이 노동시장의 균형임금은 완전경쟁 노동시장의 균형임금보다 낮다.
  \item 이 노동시장의 균형임금과 완전경쟁 노동시장의 균형임금 사이에 최저임금을 강제적으로 설정할 경우 노동의 평균요소비용과 한계요소비용이 모두 감소한다.
\end{choices}
\end{question}
\answer{5}
